\scanpage{11}

$C$ ``configurations''

$\Omega$ ``opérations'' (virtuelles)

\[\resizebox{\ifdim\width>\linewidth\linewidth\else\width\fi}{!}{$\displaystyle
 \omega\in\Omega\qquad
 \omega\longmapsto\bigl(f_\omega:C\longrightarrow
 \underset{\substack{\shortparallel\\\text{\scriptsize ``ens. des résultats \inserted{\text{virtuels}} des l'opération }\omega}}{V_\omega}\bigr)
$}\]
\margin{On pose $f_\omega(x)=\omega(x)$}

$\widetilde\Omega$ ``opérations effectuées'' :
\[
 \widetilde\Omega=\coprod_{\omega\in\Omega}V_\omega
 \xrightarrow{\pi}\Omega
\]
\[\resizebox{\ifdim\width>\linewidth\linewidth\else\width\fi}{!}{$\displaystyle
\left\{\begin{aligned}
 C(\omega,x)&=f_\omega^{-1}(\{\deleted{\unsure{f}}f_\omega(x)\})&
 \text{ombre de \inserted{\text{la conf.}} }x\in C\text{ pour l'op }\omega.\\
 &=C(\omega(x))\\
 &\text{de même on définit}\\
 C(\omega_1,\ldots\omega_\nu\,;x)&=\bigcap_{1\leq i\leq\nu}C(\omega_i\,;x)&
 \text{ombre de }x\in C\text{ pour}\\
 &&\text{la suite d'op. }\omega_1\ldots\omega_\nu
\end{aligned}\right.
$}\]

\[\resizebox{\ifdim\width>\linewidth\linewidth\else\width\fi}{!}{$\displaystyle
\begin{aligned}
 \widetilde\Omega&\xrightarrow{\widetilde\omega\mapsto C(\widetilde\omega)}
 \unsure{\mathscr P|\wp|\mathfrak P}(C)\\
 \widetilde\omega=(\omega,v)&\longmapsto f_\omega^{-1}(\{v\})
\end{aligned}
$}\]
\margin{\uline{NB} \unsure{OPS} pour la suite que les $f_\omega$ sont surj}

\[
 \widetilde\omega\longmapsto C(\widetilde\omega)
 \qquad\text{ensemble résiduel pour l'op. eff. }\widetilde\omega
\]
\[\resizebox{\ifdim\width>\linewidth\linewidth\else\width\fi}{!}{$\displaystyle
\begin{aligned}
 \widetilde\Omega^\nu&\unsure{\longrightarrow|\text{---}}\mathscr P(C)\\
 (\widetilde\omega_1,\ldots,\widetilde\omega_\nu)&\longmapsto
 C(\widetilde\omega_1,\ldots,\widetilde\omega_\nu)
 =\bigcap_{1\leq i\leq\nu}C(\widetilde\omega_i)
\end{aligned}
$}\]

ensemble résiduel pour la suite d'opérations \unsure{effectives|effectuées} $\widetilde\omega_1\ldots\widetilde\omega_\nu$.

\[\resizebox{\ifdim\width>\linewidth\linewidth\else\width\fi}{!}{$\displaystyle
\begin{aligned}
 \Omega^\nu&\ \tikz[baseline=-.2ex,x=1cm,y=1cm]{\draw[->] (0,.88)..controls(-.09,.88) and (-.10,.44)..(-.09,.30)..controls(-.09,.23) and (.04,.18)..(.48,0);}\ \text{\uline{tests} virtuels de longueur }\nu\\
 T_\nu\subset\widetilde\Omega^\nu&\quad\text{--- effectués de longueur }\nu
\end{aligned}
$}\]
\[\resizebox{\ifdim\width>\linewidth\linewidth\else\width\fi}{!}{$\displaystyle
\begin{aligned}
 C\times\Omega&\longrightarrow\widetilde\Omega\quad\deleted{\unsure{\longrightarrow|\Longleftrightarrow|\Longrightarrow}}&
 C&\longrightarrow\Gamma(\widetilde\Omega/\Omega)\\
 (x,\omega)&\longmapsto\omega(x)&
 C\times\Omega^\nu&\longrightarrow\widetilde\Omega^\nu
\end{aligned}
$}\]
\[\resizebox{\ifdim\width>\linewidth\linewidth\else\width\fi}{!}{$\displaystyle
\begin{aligned}
 C&\longrightarrow\deleted{\unsure{\Omega}}\Gamma(\widetilde\Omega^\nu/\Omega^\nu)\\
 x&\longmapsto\bigl((\omega_1,\ldots,\omega_\nu)\longmapsto
 (\omega_1(x),\ldots,\omega_\nu(x))\bigr)
\end{aligned}
$}\]

l'image de $C\times\Omega^\nu$ est appellée ens des tests \deleted{\unsure{effectués}} \inserted{\text{\unsure{effectifs|effectués}}} de lg. $\nu$. Soit $T_\nu$ l'ens de ceux-ci.
% archive sheet number: 17
